% Gliederung Stand 30.10.2026. ~5 Seiten.
\chapter{Handlungsbedarf, Anforderungen und Bewertungsmethodik}
\label{chap:handlungsbedarf}
% Kurze Kapiteleinleitung vor 5.1.

\section{Handlungsbedarf}
\label{sec:handlungsbedarf}
% Wertung der Literatur aus Kap. 3 und 4; Lücke: kein Stack verbindet gleichen Code Sim/Real,
% Laienbedienung, Sicherheit und Erweiterbarkeit. Kandidaten: Yu2026 (G1-Studie ohne Realversuch), Tong2024 (sicherer Sim-Real-Übergang offen).

\section{Anforderungen}
\label{sec:anforderungen}
% Tabelle A1–A7 mit Kennzahlen (Startwerte, Herleitung hier begründen):
% A1 aufgabenunabhängig (F4) | A2 gleicher Code Sim/Real (F1) | A3 echtzeitfähig auf Standard-PC (F2)
% A4 realitätsnahe Simulation (F2) | A5 sicher im Realbetrieb (F3) | A6 bedienbar ohne Vorerfahrung (F3)
% A7 erweiterbar (F4)

\section{Bewertungsmethodik}
\label{sec:bewertungsmethodik}

\subsection{Zuordnung von Kennzahlen zu Anforderungen und Forschungsfragen}
\label{sec:kennzahlen}

\subsection{Versuchsplan}
\label{sec:versuchsplan}
% Simulation und Realbetrieb, gleiche Kennzahlen wo möglich.

\subsection{Testumgebung}
\label{sec:testumgebung}
% Referenz-PC (Hardware, Softwareversionen), automatisierte Tests, Fake-G1.
